%!TEX root = ../template.tex
%%%%%%%%%%%%%%%%%%%%%%%%%%%%%%%%%%%%%%%%%%%%%%%%%%%%%%%%%%%%%%%%%%%%
%% chapter7.tex
%% NOVA thesis document file
%%
%% Chapter with lots of dummy text
%%%%%%%%%%%%%%%%%%%%%%%%%%%%%%%%%%%%%%%%%%%%%%%%%%%%%%%%%%%%%%%%%%%%

\typeout{NT FILE chapter7.tex}%

\chapter{Unit Testing}
\label{cha:UnitTesting}

This chapter presents the unit tests used to verify the correctness and robustness of the attribute-grammar implementation. 
The tests cover core scenarios—successful evaluation of synthesized and inherited attributes, deterministic left-to-right child evaluation, 
and conformance between parse trees and production rules—as well as boundary cases that exercise error handling, such as type mismatches, 
division by zero, invalid child occurrences and dependency cycles.

\section{The Tests}

\subsection{\texttt{test\_parseTree} (Grammar \texttt{ag2}, word \texttt{3+2+9\textasciitilde})}
Parse the word \texttt{3+2+9\textasciitilde} using the CFG projected from \texttt{ag2}, convert the CFG tree to an AG tree, and evaluate attributes with \texttt{calcAttributes}. The shape features a head \texttt{S$\rightarrow$E}, with \texttt{E$\rightarrow$F\ X} and a right-branching chain of \texttt{X} nodes for the two \texttt{+} and the final \texttt{\textasciitilde}.


\vspace{1cm}
\begin{tikzpicture}[
  level 1/.style={sibling distance=48mm},
  level 2/.style={sibling distance=30mm},
  level 3/.style={sibling distance=22mm},
  level 4/.style={sibling distance=18mm},
  edge from parent/.style={draw, -latex},
  every node/.style={draw, ellipse, align=center}
]
\node {S}
  child {node {E}
    child {node {F}
      child {node {3}}
    }
    child {node {X}
      child {node {+}}
      child {node {F}
        child {node {2}}
      }
      child {node {X}
        child {node {+}}
        child {node {F}
          child {node {9}}
        }
        child {node {X}
          child {node {\textasciitilde}}
        }
      }
    }
  };
\end{tikzpicture}

\vspace{1cm}

\subsection{\texttt{test\_ag1\_simple\_synthesized} (\texttt{ag1}, tree \texttt{pt1})}
Validate a simple synthesized attribute \texttt{v} over multiplication: \texttt{S$\rightarrow$E}, \texttt{E$\rightarrow$E\ *\ F}, inner \texttt{E$\rightarrow$F}, and leaves \texttt{F$\rightarrow$3}, \texttt{F$\rightarrow$2}. Expected result after evaluation: \texttt{v(S)=6}.%
% 

\vspace{1cm}
\begin{tikzpicture}[
  level 1/.style={sibling distance=45mm},
  level 2/.style={sibling distance=30mm},
  level 3/.style={sibling distance=20mm},
  edge from parent/.style={draw, -latex},
  every node/.style={draw, ellipse, align=center}
]
\node {S}
  child {node {E}
    child {node {E}
      child {node {F}
        child {node {3}}
      }
    }
    child {node {\*}}
    child {node {F}
      child {node {2}}
    }
  };
\end{tikzpicture}

\vspace{1cm}

\subsection{\texttt{test\_ag2\_simple\_synthesized\_and\_inherited} (\texttt{ag2}, tree \texttt{pt2})}
 Check inherited flow \texttt{d} across the \texttt{X} chain and synthesized result \texttt{r}/\texttt{v}. The fixture \texttt{pt2} encodes the same structure as the parsed word \texttt{3+2+9\textasciitilde}. Expected: \texttt{v(S)=14}.%
% 

\vspace{1cm}
\begin{tikzpicture}[
  level 1/.style={sibling distance=48mm},
  level 2/.style={sibling distance=30mm},
  level 3/.style={sibling distance=22mm},
  level 4/.style={sibling distance=18mm},
  edge from parent/.style={draw, -latex},
  every node/.style={draw, ellipse, align=center}
]
\node {S}
  child {node {E}
    child {node {F}
      child {node {3}}
    }
    child {node {X}
      child {node {+}}
      child {node {F}
        child {node {2}}
      }
      child {node {X}
        child {node {+}}
        child {node {F}
          child {node {9}}
        }
        child {node {X}
          child {node {\textasciitilde}}
        }
      }
    }
  };
\end{tikzpicture}

\vspace{1cm}

\subsection{\texttt{test\_ag3\_simpler\_synthesized\_and\_inherited} (\texttt{ag3}, tree \texttt{pt3})}
Minimal inherited and synthesized flow: \texttt{S$\rightarrow$E \{d(E)=5, v(S)=v(E)\}}, \texttt{E$\rightarrow$\textasciitilde \{v(E)=d(E)+1\}}. Expected: \texttt{v(S)=6}.%
% 

\vspace{1cm}
\begin{tikzpicture}[
  level 1/.style={sibling distance=35mm},
  edge from parent/.style={draw, -latex},
  every node/.style={draw, ellipse, align=center}
]
\node {S}
  child {node {E}
    child {node {\textasciitilde}}
  };
\end{tikzpicture}

\vspace{1cm}

\subsection{\texttt{test\_ag7\_default\_index} (\texttt{ag7}, tree \texttt{pt7})}
Confirm default child indices and explicit \texttt{C1}/\texttt{C2} occurrences on \texttt{S$\rightarrow$C\ C} are handled correctly when applying inherited assignments. Leaves are \texttt{n} and \texttt{m}.%
% 

\vspace{1cm}
\begin{tikzpicture}[
  level 1/.style={sibling distance=40mm},
  level 2/.style={sibling distance=24mm},
  edge from parent/.style={draw, -latex},
  every node/.style={draw, ellipse, align=center}
]
\node {S}
  child {node {C}
    child {node {n}}
  }
  child {node {C}
    child {node {m}}
  };
\end{tikzpicture}

\vspace{1cm}

\subsection{\texttt{test\_ag8\_boolean\_flags} (\texttt{ag8}, tree \texttt{pt8})}
Evaluate a boolean synthesized attribute \texttt{b} over terminals \texttt{x} and \texttt{y}. Fixture expands \texttt{S$\rightarrow$A\ B} with \texttt{A$\rightarrow$x} and \texttt{B$\rightarrow$y}.%
% 

\vspace{1cm}
\begin{tikzpicture}[
  level 1/.style={sibling distance=40mm},
  level 2/.style={sibling distance=24mm},
  edge from parent/.style={draw, -latex},
  every node/.style={draw, ellipse, align=center}
]
\node {S}
  child {node {A}
    child {node {x}}
  }
  child {node {B}
    child {node {y}}
  };
\end{tikzpicture}

\vspace{1cm}

\subsection{\texttt{test\_ag9\_list\_length} (\texttt{ag9}, tree \texttt{pt9 = make\_list 30})}
Compute list length using a synthesized attribute \texttt{v}. The builder \texttt{make\_list} nests \texttt{L} to the left: for \texttt{n=1}, \texttt{L$\rightarrow$I}; for \texttt{n>1}, \texttt{L$\rightarrow$L , I}. The actual test uses \texttt{n=30}. Below is an illustrative tree for \texttt{n=4}; the \texttt{n=30} instance continues the same left-nested pattern. Expected at the root: length \texttt{30}.

\vspace{1cm}
\begin{tikzpicture}[
  level 1/.style={sibling distance=55mm},
  level 2/.style={sibling distance=40mm},
  level 3/.style={sibling distance=30mm},
  level 4/.style={sibling distance=24mm},
  edge from parent/.style={draw, -latex},
  every node/.style={draw, ellipse, align=center}
]
\node {L}
  child {node {L}
    child {node {L}
      child {node {I}
        child {node {a}}
      }
    }
    child {node {,}}
    child {node {I}
      child {node {a}}
    }
  }
  child {node {,}}
  child {node {I}
    child {node {a}}
  };
\end{tikzpicture}

\vspace{1cm}

\subsection{\texttt{test\_ag10\_inherited\_default} (\texttt{ag10}, tree \texttt{pt10})}
Apply default inherited values to both children of \texttt{S$\rightarrow$C\ F}, then synthesize \texttt{v} at \texttt{S}. Leaves are \texttt{n} and \texttt{m}. Expected: \texttt{v(S)=3}.%
% 

\vspace{1cm}
\begin{tikzpicture}[
  level 1/.style={sibling distance=40mm},
  level 2/.style={sibling distance=24mm},
  edge from parent/.style={draw, -latex},
  every node/.style={draw, ellipse, align=center}
]
\node {S}
  child {node {C}
    child {node {n}}
  }
  child {node {F}
    child {node {m}}
  };
\end{tikzpicture}

\vspace{1cm}

\subsection{\texttt{test\_ag\_expr\_ext} (\texttt{ag\_expr\_ext}, tree \texttt{pt\_expr\_ext})}
Extended arithmetic with \texttt{+}, \texttt{-}, \texttt{\*}, \texttt{/}, and parentheses. Attributes: value \texttt{v}, pretty string \texttt{s}, and height \texttt{h}. Fixture encodes \texttt{3 + 2 * (9 - 5) / 2 - 1} under nonterminals \texttt{E/T/F/N}. Expected at the root: \texttt{v(S)=6}, \texttt{s(S)=``3+2*(9-5)/2-1''}.%
% 

\vspace{1cm}
\begin{tikzpicture}[
  level 1/.style={sibling distance=60mm},
  level 2/.style={sibling distance=45mm},
  level 3/.style={sibling distance=32mm},
  level 4/.style={sibling distance=26mm},
  level 5/.style={sibling distance=22mm},
  edge from parent/.style={draw, -latex},
  every node/.style={draw, ellipse, align=center}
]
\node {S}
  child {node {E}
    child {node {E}
      child {node {E}
        child {node {T}
          child {node {F}
            child {node {N}
              child {node {3}}
            }
          }
        }
      }
      child {node {+}}
      child {node {T}
        child {node {T}
          child {node {T}
            child {node {F}
              child {node {N}
                child {node {2}}
              }
            }
          }
          child {node {\*}}
          child {node {F}
            child {node {(}}
            child {node {E}
              child {node {E}
                child {node {T}
                  child {node {F}
                    child {node {N}
                      child {node {9}}
                    }
                  }
                }
              }
              child {node {-}}
              child {node {T}
                child {node {F}
                  child {node {N}
                    child {node {5}}
                  }
                }
              }
            }
            child {node {)}}
          }
        }
        child {node {/}}
        child {node {F}
          child {node {N}
            child {node {2}}
          }
        }
      }
    }
    child {node {-}}
    child {node {T}
      child {node {F}
        child {node {N}
          child {node {1}}
        }
      }
    }
  };
\end{tikzpicture}

\vspace{1cm}